% ==============================================================================
% CAPÍTULO 02 - ESTADO DEL ARTE
% ==============================================================================

\chapter{Estado del Arte}
\label{cap:estado_arte}
\justifying

% ==============================================================================
% INTRODUCCIÓN
% ==============================================================================


\begin{flushleft}
\justifying

La transformación digital ha impulsado el desarrollo de plataformas orientadas a la gestión electrónica de denuncias, quejas y mecanismos de protección de derechos. En este contexto, las organizaciones gubernamentales, corporativas y académicas han incorporado sistemas digitales capaces de automatizar la recepción de reportes, el seguimiento de expedientes y la administración de evidencia electrónica.

Sin embargo, la evolución de estas plataformas no ha sido homogénea. Mientras algunos sistemas priorizan la validez jurídica y el cumplimiento normativo, otros se enfocan en la anonimización de datos, la movilidad, la trazabilidad o la integración de mecanismos de seguridad informática. Esto ha provocado que muchas soluciones resulten insuficientes para contextos institucionales complejos, particularmente en entornos universitarios donde convergen aspectos administrativos, jurídicos, académicos y de protección de derechos humanos.

En el caso de la DDP-IPN, la necesidad de una plataforma especializada implica considerar elementos que van más allá de un sistema convencional de recepción de denuncias o quejas. Entre los aspectos críticos identificados destacan:

\begin{itemize}

    \item Protección e integridad de la información sensible.
    
    \item Automatización de procesos administrativos.
    
    \item Escalabilidad de la infraestructura tecnológica.
    
    \item Trazabilidad de expedientes digitales.
    
    \item Protección de identidad y anonimización del denunciante.
    
    \item Integración institucional entre distintas áreas administrativas.
    
    \item Adaptabilidad a procesos universitarios complejos.

\end{itemize}

En este sentido, el presente capítulo analiza diversas plataformas nacionales e internacionales relacionadas con sistemas de denuncia digital y gestión institucional, con el objetivo de identificar las capacidades tecnológicas existentes, las limitaciones arquitectónicas predominantes y las brechas funcionales que justifican el desarrollo de una solución especializada para el Instituto Politécnico Nacional.

\end{flushleft}

% ==============================================================================
% CRITERIOS DE EVALUACIÓN
% ==============================================================================

\section{Criterios de Evaluación}

\begin{flushleft}
\justifying

Con el propósito de realizar un análisis técnico homogéneo entre las plataformas seleccionadas, se definieron criterios de evaluación orientados tanto a aspectos funcionales como arquitectónicos.

Los principales criterios considerados son los siguientes:

\begin{itemize}

    \item \textbf{Escalabilidad Arquitectónica:}
    Capacidad del sistema para soportar crecimiento en usuarios, servicios y procesamiento sin degradar significativamente su rendimiento.

    \item \textbf{Seguridad y Confidencialidad:}
    Mecanismos implementados para proteger la información sensible, incluyendo cifrado, autenticación y control de acceso.

    \item \textbf{Anonimización y Protección de Identidad:}
    Nivel de protección ofrecido al denunciante para preservar su identidad y confidencialidad.

    \item \textbf{Trazabilidad y Seguimiento:}
    Capacidad de monitorear el estado y avance de expedientes o denuncias digitales.

    \item \textbf{Automatización Administrativa:}
    Uso de tecnologías orientadas a reducir procesos manuales mediante reglas de negocio, clasificación automática o asistencia inteligente.

    \item \textbf{Soporte Multimedia:}
    Capacidad para almacenar y gestionar evidencia digital como imágenes, documentos, audio y video.

    \item \textbf{Interoperabilidad Institucional:}
    Facilidad para integrarse con otros servicios, sistemas o plataformas institucionales.

    \item \textbf{Adaptación al Entorno Universitario:}
    Nivel de adecuación funcional y administrativa para instituciones educativas públicas.

\end{itemize}

A partir de estos criterios se desarrolló el análisis comparativo de las plataformas seleccionadas.

\end{flushleft}

% ==============================================================================
% PLATAFORMAS GUBERNAMENTALES
% ==============================================================================

\section{Plataformas Gubernamentales}

En esta sección se analiza el panorama actual de las herramientas tecnológicas implementadas por el Estado mexicano para la gestión, recepción y seguimiento de denuncias ciudadanas. A través de la revisión de soluciones del ámbito federal, como el Sistema Integral de Denuncias Ciudadanas (SIDEC), y de iniciativas locales, como la plataforma Denuncia Digital de la Ciudad de México, se examinan sus capacidades funcionales, su arquitectura tecnológica y los retos pendientes en materia de interoperabilidad, automatización y descentralización.

\subsection{Sistema Integral de Denuncias Ciudadanas (SIDEC)}

\begin{flushleft}
\justifying

El Sistema Integral de Denuncias Ciudadanas (SIDEC) constituye una de las principales plataformas gubernamentales de recepción de quejas administrativas en México, centralizando denuncias relacionadas con servidores públicos federales \cite{6}.

Desde una perspectiva funcional, el sistema permite la recepción continua de denuncias mediante formularios electrónicos y proporciona mecanismos básicos de seguimiento mediante folios únicos y contraseñas de acceso \cite{7}. Asimismo, incorpora soporte para evidencia multimedia, incluyendo documentos, audio y video.

No obstante, el análisis arquitectónico del SIDEC evidencia diversas limitaciones técnicas relevantes:

\begin{itemize}

    \item Arquitectura predominantemente centralizada.
    
    \item Procesos administrativos parcialmente manuales.
    
    \item Seguimiento limitado para el denunciante.
    
    \item Ausencia de automatización inteligente.
    
    \item Escasa modularidad de servicios.

\end{itemize}

\end{flushleft}

% ------------------------------------------------------------------------------
% DENUNCIA DIGITAL
% ------------------------------------------------------------------------------

\subsection{Denuncia Digital y Sistemas Estatales}

\begin{flushleft}
\justifying

Diversas entidades federativas han incorporado plataformas digitales orientadas a la recepción de denuncias y querellas. Entre los casos más representativos se encuentra el sistema \textbf{Denuncia Digital} implementado por la Ciudad de México \cite{9}.

Esta plataforma permite iniciar procedimientos relacionados con delitos no violentos mediante autenticación electrónica utilizando mecanismos como la \textit{Llave CDMX} o firma electrónica gubernamental.

El análisis comparativo entre entidades federativas muestra una importante fragmentación tecnológica, donde algunas plataformas permiten videodenuncias y ratificación remota, mientras otras continúan dependiendo de modelos híbridos o procesos presenciales.

A pesar de representar un avance significativo en digitalización administrativa, estas soluciones presentan limitaciones relacionadas con interoperabilidad, automatización y protección avanzada de identidad.

\end{flushleft}

% ------------------------------------------------------------------------------
% MP VIRTUAL
% ------------------------------------------------------------------------------

\subsection{Servicios de Procuración Digital: MP Virtual y Denuncia Anónima}

\begin{flushleft}
\justifying

La Fiscalía General de Justicia de la Ciudad de México ha incorporado servicios digitales complementarios como el \textbf{MP Virtual} y el sistema de \textbf{Denuncia Anónima} \cite{11}.

El módulo MP Virtual permite formalizar determinados procedimientos jurídicos de manera remota, facilitando la validación documental y reduciendo la necesidad de comparecencia presencial \cite{12}.

Por otra parte, el sistema de denuncia anónima implementa mecanismos básicos de protección de identidad para reportes relacionados con delitos de alto impacto.

Sin embargo, estas plataformas continúan dependiendo de infraestructuras centralizadas y procesos administrativos tradicionales, además de carecer de componentes avanzados de inteligencia artificial, clasificación automática y analítica institucional.

\end{flushleft}

% ==============================================================================
% APLICACIONES MÓVILES Y CIUDADANAS
% ==============================================================================

\section{Aplicaciones Móviles y Plataformas Ciudadanas}

% ------------------------------------------------------------------------------
% INCORRUPTIBLE
% ------------------------------------------------------------------------------

\subsection{Aplicaciones de Denuncia Georreferenciada}

\begin{flushleft}
\justifying

La evolución de los servicios digitales ha impulsado el desarrollo de aplicaciones móviles orientadas a la denuncia ciudadana y participación colaborativa.

Uno de los referentes más representativos es la aplicación \textbf{Incorruptible}, diseñada para registrar actos de corrupción mediante mecanismos de geolocalización y participación ciudadana \cite{13}.

Entre las principales funcionalidades identificadas destacan:

\begin{itemize}

    \item Registro georreferenciado de incidentes.
    
    \item Publicación colaborativa de reportes.
    
    \item Integración de evidencia multimedia.
    
    \item Cuantificación económica de daños reportados.

\end{itemize}

De forma complementaria, el \textbf{Sistema de Atención Mexiquense (SAM)} integra servicios web, móviles y telefónicos para canalizar denuncias hacia distintas dependencias administrativas \cite{14}.

A pesar de su accesibilidad y facilidad de uso, este tipo de aplicaciones presentan limitaciones importantes relacionadas con trazabilidad jurídica, validación institucional y protección de datos sensibles.

\end{flushleft}

% ==============================================================================
% PLATAFORMAS INTERNACIONALES
% ==============================================================================

\section{Plataformas Internacionales}

\begin{flushleft}
\justifying

El mercado internacional de plataformas de denuncia y cumplimiento normativo ha desarrollado soluciones avanzadas enfocadas principalmente en sectores corporativos y regulatorios.

Estas plataformas destacan por incorporar mecanismos avanzados de seguridad, anonimización y administración de casos; sin embargo, suelen presentar limitaciones relacionadas con costos de implementación, dependencia de proveedores externos y adaptación a marcos institucionales específicos.

\end{flushleft}

% ------------------------------------------------------------------------------
% WHISTLELINK
% ------------------------------------------------------------------------------

\subsection{Whistlelink}

\begin{flushleft}
\justifying

Whistlelink es una plataforma orientada a la gestión anónima de denuncias corporativas, destacando por su soporte multilingüe y canales bidireccionales de comunicación anónima \cite{15}.

Su principal ventaja radica en la reducción de barreras de acceso para denunciantes internacionales; sin embargo, presenta limitaciones relacionadas con personalización institucional avanzada y adaptación a procesos administrativos complejos.

\end{flushleft}

% ------------------------------------------------------------------------------
% EQS
% ------------------------------------------------------------------------------

\subsection{EQS Integrity Line}

\begin{flushleft}
\justifying

EQS Integrity Line constituye una de las soluciones corporativas más robustas del mercado europeo, diseñada para cumplir con normativas internacionales de protección al denunciante \cite{16}.

La plataforma incorpora herramientas avanzadas de administración de casos, monitoreo organizacional y generación de reportes.

No obstante, su complejidad operativa y elevada carga administrativa pueden dificultar su implementación en instituciones educativas públicas.

\end{flushleft}

% ------------------------------------------------------------------------------
% NAVEX
% ------------------------------------------------------------------------------

\subsection{NAVEX / WhistleB}

\begin{flushleft}
\justifying

NAVEX y su solución WhistleB utilizan infraestructura en la nube para garantizar alta disponibilidad y escalabilidad \cite{17}.

Entre sus fortalezas destacan:

\begin{itemize}

    \item Integración con plataformas GRC.
    
    \item Infraestructura escalable.
    
    \item Gestión centralizada de cumplimiento.
    
    \item Administración avanzada de casos.

\end{itemize}

Sin embargo, la dependencia de proveedores externos puede representar restricciones relacionadas con soberanía de datos y políticas institucionales gubernamentales.

\end{flushleft}

% ------------------------------------------------------------------------------
% WHISTLEBLOWER SOFTWARE
% ------------------------------------------------------------------------------

\subsection{Whistleblower Software}

\begin{flushleft}
\justifying

Whistleblower Software destaca por implementar mecanismos avanzados de cifrado y modelos de arquitectura Zero-Knowledge \cite{18}.

Este enfoque fortalece significativamente la confidencialidad de la información y la protección del denunciante.

A pesar de ello, su orientación hacia entornos corporativos limita su adecuación directa a estructuras universitarias públicas.

\end{flushleft}

% ==============================================================================
% PLATAFORMAS ACADÉMICAS
% ==============================================================================

\section{Plataformas Académicas e Institucionales}

En esta sección se examinan las principales herramientas tecnológicas e iniciativas digitales desarrolladas por instituciones educativas y organismos públicos para la atención de incidencias en el ámbito académico. A través del análisis de soluciones específicas como BullyButton, Ventanilla Escolar, Alerta IPN y el sistema Denuncia Línea de la UNAM, se evalúan sus alcances operativos y protocolos de respuesta. Asimismo, se identifican las brechas tecnológicas comunes en estos entornos, orientadas principalmente a las limitaciones en la gestión jurídica de expedientes, la automatización del triaje y la falta de mecanismos avanzados para la protección de la identidad de los usuarios.

% ------------------------------------------------------------------------------
% BULLYBUTTON
% ------------------------------------------------------------------------------

\subsection{BullyButton}

\begin{flushleft}
\justifying

BullyButton es una plataforma enfocada en la detección y reporte de casos de acoso escolar \cite{21}.

Su diseño prioriza la simplicidad operativa y facilidad de uso para estudiantes; sin embargo, carece de mecanismos formales de trazabilidad jurídica y administración avanzada de expedientes.

\end{flushleft}

% ------------------------------------------------------------------------------
% VENTANILLA ESCOLAR
% ------------------------------------------------------------------------------

\subsection{Ventanilla Escolar}

\begin{flushleft}
\justifying

Ventanilla Escolar funciona como un canal de orientación y comunicación entre instituciones educativas y autoridades federales \cite{22}.

Aunque centraliza información relevante para usuarios educativos, no incorpora funcionalidades completas relacionadas con seguimiento de denuncias y administración formal de expedientes.

\end{flushleft}

% ------------------------------------------------------------------------------
% ALERTA IPN
% ------------------------------------------------------------------------------

\subsection{Alerta IPN}

\begin{flushleft}
\justifying

Alerta IPN constituye uno de los principales mecanismos institucionales de reporte dentro del Instituto Politécnico Nacional \cite{23}.

Su funcionalidad se orienta principalmente a reportes de emergencia y situaciones de riesgo dentro de instalaciones institucionales.

No obstante, la plataforma presenta limitaciones relacionadas con:

\begin{itemize}

    \item Gestión jurídica de expedientes.
    
    \item Clasificación administrativa automatizada.
    
    \item Seguimiento integral de casos.
    
    \item Protección avanzada de identidad.
    
    \item Automatización inteligente de procesos.

\end{itemize}

\end{flushleft}

% ------------------------------------------------------------------------------
% UNAM
% ------------------------------------------------------------------------------

\subsection{UNAM - Denuncia Línea}

\begin{flushleft}
\justifying

La plataforma Denuncia Línea de la UNAM representa uno de los referentes más relevantes en gestión universitaria de denuncias en México \cite{24}.

Su principal fortaleza radica en la existencia de protocolos institucionales consolidados y procedimientos administrativos claramente definidos.

Sin embargo, el análisis tecnológico evidencia limitaciones relacionadas con:

\begin{itemize}

    \item Escalabilidad arquitectónica.
    
    \item Integración de inteligencia artificial.
    
    \item Automatización de triaje.
    
    \item Modularidad de servicios.
    
    \item Gestión inteligente de expedientes.

\end{itemize}

\end{flushleft}

% ==============================================================================
% TABLA COMPARATIVA GENERAL
% ==============================================================================

\section{Tabla Comparativa General}

\begin{flushleft}
\justifying

Con base en los criterios de evaluación previamente definidos, se realizó una comparación general de las plataformas analizadas con el objetivo de identificar sus capacidades funcionales, limitaciones arquitectónicas y nivel de adaptación a entornos institucionales complejos.

La Tabla \ref{tab:comparativa_general} resume las principales características identificadas durante la investigación.

\end{flushleft}

\begin{table}[H]
\centering
\caption{Comparativa general de plataformas de denuncia digital}
\label{tab:comparativa_general}

\renewcommand{\arraystretch}{1.4}

\resizebox{\textwidth}{!}{

\begin{tabular}{|p{3.2cm}|c|c|c|c|c|c|c|}
\hline

\textbf{Plataforma} &
\textbf{Seguridad} &
\textbf{Anonimato} &
\textbf{Seguimiento} &
\textbf{Multimedia} &
\textbf{Automatización} &
\textbf{Escalabilidad} &
\textbf{Adaptación Universitaria}
\\

\hline

SIDEC &
Media &
Baja &
Media &
Sí &
Baja &
Media &
No
\\

\hline

Denuncia Digital CDMX &
Media &
Baja &
Media &
Sí &
Baja &
Media &
No
\\

\hline

MP Virtual &
Media &
Media &
Media &
Sí &
Baja &
Media &
No
\\

\hline

Incorruptible &
Baja &
Media &
Baja &
Sí &
Baja &
Alta &
No
\\

\hline

SAM &
Media &
Baja &
Media &
Sí &
Baja &
Media &
No
\\

\hline

Whistlelink &
Alta &
Alta &
Alta &
Sí &
Media &
Alta &
Parcial
\\

\hline

EQS Integrity Line &
Alta &
Alta &
Alta &
Sí &
Alta &
Alta &
Parcial
\\

\hline

NAVEX / WhistleB &
Alta &
Alta &
Alta &
Sí &
Alta &
Alta &
Parcial
\\

\hline

Whistleblower Software &
Muy Alta &
Muy Alta &
Alta &
Sí &
Alta &
Alta &
Parcial
\\

\hline

BullyButton &
Media &
Media &
Baja &
Limitado &
Baja &
Media &
Sí
\\

\hline

Ventanilla Escolar &
Media &
Baja &
Baja &
No &
Baja &
Media &
Sí
\\

\hline

Alerta IPN &
Media &
Baja &
Media &
Sí &
Baja &
Media &
Sí
\\

\hline

UNAM Denuncia Línea &
Media &
Media &
Alta &
Sí &
Media &
Media &
Sí
\\

\hline

\end{tabular}

}

\end{table}

% ==============================================================================
% BRECHAS TECNOLÓGICAS
% ==============================================================================

\section{Brechas Tecnológicas Identificadas}

\begin{flushleft}
\justifying

El análisis comparativo permitió identificar que, aunque existen múltiples plataformas orientadas a la gestión digital de denuncias y quejas, ninguna integra de manera completa los requerimientos tecnológicos, administrativos y de protección institucional necesarios para un entorno universitario público como el Instituto Politécnico Nacional.

Entre las principales brechas identificadas destacan:

\begin{itemize}

    \item Dependencia de arquitecturas centralizadas tradicionales.
    
    \item Escasa automatización inteligente de procesos administrativos.
    
    \item Limitados mecanismos avanzados de anonimización y protección de identidad.
    
    \item Baja interoperabilidad entre servicios institucionales.
    
    \item Deficiencias en trazabilidad integral de expedientes digitales.
    
    \item Escasa adaptación a estructuras universitarias complejas.
    
    \item Ausencia de arquitecturas híbridas escalables orientadas a instituciones académicas públicas.

\end{itemize}

Estas limitaciones justifican la necesidad de desarrollar una plataforma especializada que combine seguridad informática avanzada, modularidad arquitectónica, automatización administrativa y adaptación específica al contexto institucional de la DDP-IPN.

\end{flushleft}